%% ============================================================
%% generative.tex — 第10章 拡散モデルと生成モデル
%%
%% 本章では，VAE，拡散モデル（DDPM），
%% およびPhysics-Informed Neural Networks（PINN）を解説する．
%% ============================================================

\chapter{拡散モデルと生成モデル}
\label{chap:generative}

本章では，データの分布を学習し新たなサンプルを生成する
\textbf{生成モデル（Generative Models）}について学ぶ．
特に，変分オートエンコーダ（VAE）と拡散モデル（Diffusion Models）を中心に解説し，
物理法則を組み込んだニューラルネットワーク（PINN）についても触れる．
これらの技術は，交通シミュレーション，軌跡生成，都市計画など
多様な応用可能性を持つ．

% ------------------------------------------------------------------
\section{生成モデルの確率的フレームワーク}
\label{sec:generative-framework}
% ------------------------------------------------------------------

\subsection{生成モデルの目標}

生成モデルの根本的な目標は，\textbf{未知のデータ分布 $p_{\mathrm{data}}(\boldsymbol{x})$ を学習し，
それに近い分布 $p_\theta(\boldsymbol{x})$ を構築する}ことである．

識別モデル $p(y|\boldsymbol{x})$ が「与えられた入力から出力を予測する」のに対し，
生成モデル $p(\boldsymbol{x})$ は「データそのものの本質的な構造を理解し，
新たなサンプルを創造できる」という意味で，より深い学習を実現している．

\subsection{潜在変数モデル}

実際のデータは多くの場合，\textbf{低次元の潜在構造}を持っている．
この直感を形式化するのが\textbf{潜在変数モデル}である：
\begin{equation}
  p_\theta(\boldsymbol{x}) = \int p_\theta(\boldsymbol{x}|\boldsymbol{z}) p(\boldsymbol{z}) d\boldsymbol{z}.
  \label{eq:latent-model}
\end{equation}

ここで：
\begin{itemize}
  \item $\boldsymbol{z} \in \mathbb{R}^L$：潜在変数（通常 $L < D$）
  \item $p(\boldsymbol{z})$：事前分布（通常は $\mathcal{N}(\mathbf{0}, \mathbf{I})$）
  \item $p_\theta(\boldsymbol{x}|\boldsymbol{z})$：生成分布
\end{itemize}

\subsection{周辺尤度の計算困難性}

最大尤度推定 $\theta^* = \arg\max_\theta \sum_{i=1}^N \log p_\theta(\boldsymbol{x}_i)$ を
行いたいが，式\eqref{eq:latent-model}の積分は一般に解析的に計算できない．

この困難に対して，複数のアプローチが存在する：

\begin{table}[hbtp]
  \centering
  \caption{生成モデルのアプローチ比較．}
  \label{tab:generative-approaches}
  \begin{tabular}{lll}
    \toprule
    アプローチ & 長所 & 短所 \\
    \midrule
    VAE & 訓練が安定，解釈可能 & 出力がぼけやすい \\
    GAN & 高品質な生成 & 訓練が不安定，mode collapse \\
    拡散モデル & 高品質・多様，安定 & 生成が遅い \\
    フローモデル & 正確な尤度評価 & 構造が複雑 \\
    \bottomrule
  \end{tabular}
\end{table}

% ------------------------------------------------------------------
\section{変分推論と VAE}
\label{sec:vae}
% ------------------------------------------------------------------

\subsection{変分推論の基本思想}

与えられたデータ $\boldsymbol{x}$ に対する事後分布 $p_\theta(\boldsymbol{z}|\boldsymbol{x})$ を
直接計算することは困難である．
\textbf{変分推論}では，計算しやすい確率分布族 $q_\phi(\boldsymbol{z}|\boldsymbol{x})$ で
真の事後分布を近似する．

\subsection{ELBO（Evidence Lower Bound）}

対数周辺尤度に対する下限（ELBO）が導出できる：
\begin{equation}
  \log p_\theta(\boldsymbol{x}) \geq \mathcal{L}(\theta, \phi; \boldsymbol{x}) = \mathbb{E}_{q_\phi(\boldsymbol{z}|\boldsymbol{x})}\left[\log p_\theta(\boldsymbol{x}|\boldsymbol{z})\right] - D_{\mathrm{KL}}(q_\phi(\boldsymbol{z}|\boldsymbol{x}) \| p(\boldsymbol{z})).
  \label{eq:elbo}
\end{equation}

ELBOの2つの項は以下のように解釈される：
\begin{itemize}
  \item \textbf{再構成誤差}：$\boldsymbol{z}$ から $\boldsymbol{x}$ をどの程度正確に再構成できるか
  \item \textbf{KL正則化}：学習された事後分布が事前分布からどの程度乖離しているか
\end{itemize}

\subsection{VAE のアーキテクチャ}

\textbf{変分オートエンコーダ（VAE）}\cite{Kingma2014}は，
変分推論をニューラルネットワークで実装したものである．

\textbf{エンコーダ} $q_\phi(\boldsymbol{z}|\boldsymbol{x})$：
\begin{equation}
  q_\phi(\boldsymbol{z}|\boldsymbol{x}) = \mathcal{N}(\boldsymbol{\mu}_\phi(\boldsymbol{x}), \mathrm{diag}(\boldsymbol{\sigma}_\phi(\boldsymbol{x})^2)).
\end{equation}

\textbf{デコーダ} $p_\theta(\boldsymbol{x}|\boldsymbol{z})$：
潜在変数 $\boldsymbol{z}$ から元のデータ空間への写像．

\subsection{再パラメータ化トリック}

サンプリングは微分不可能であるため，\textbf{再パラメータ化トリック}を用いる：
\begin{equation}
  \boldsymbol{z} = \boldsymbol{\mu}_\phi(\boldsymbol{x}) + \boldsymbol{\sigma}_\phi(\boldsymbol{x}) \odot \boldsymbol{\epsilon}, \quad \boldsymbol{\epsilon} \sim \mathcal{N}(\mathbf{0}, \mathbf{I}).
  \label{eq:reparameterization}
\end{equation}

これにより，$\phi$ に関する勾配を計算でき，標準的な自動微分で実装可能となる．

\subsection{KL 項の解析解}

ガウス分布の場合，KL項は解析的に計算できる：
\begin{equation}
  D_{\mathrm{KL}}(q_\phi \| p) = \frac{1}{2} \sum_{j=1}^L \left(\mu_j^2 + \sigma_j^2 - \log \sigma_j^2 - 1\right).
  \label{eq:kl-gaussian}
\end{equation}

\subsection{VAE の限界}

VAEには以下の既知の問題がある：
\begin{enumerate}
  \item \textbf{出力のぼけ}：再構成誤差が平均的な出力を好む傾向
  \item \textbf{事後崩壊}：KL項がゼロに近づき，潜在変数が無視される現象
\end{enumerate}

% ------------------------------------------------------------------
\section{拡散モデル（Diffusion Models）}
\label{sec:diffusion}
% ------------------------------------------------------------------

\subsection{スコアマッチング}

拡散モデルの理論的基盤として，\textbf{スコア関数}の概念が重要である：
\begin{equation}
  \mathbf{s}_\theta(\boldsymbol{x}) = \nabla_{\boldsymbol{x}} \log p_\theta(\boldsymbol{x}).
  \label{eq:score}
\end{equation}

スコアマッチングの目標は，モデルのスコア関数をデータのスコア関数に合わせることである．

\subsection{DDPM（Denoising Diffusion Probabilistic Models）}

\textbf{DDPM}\cite{Ho2020}は，データを段階的にノイズで汚す過程と，
その逆過程を学習する．

\subsubsection{前向き過程（拡散過程）}

データ $\boldsymbol{x}_0$ を段階的にノイズで汚す：
\begin{equation}
  q(\boldsymbol{x}_t | \boldsymbol{x}_{t-1}) = \mathcal{N}(\boldsymbol{x}_t; \sqrt{1-\beta_t}\boldsymbol{x}_{t-1}, \beta_t \mathbf{I}),
  \label{eq:forward-process}
\end{equation}
ここで $\beta_1, \ldots, \beta_T$ は固定されたノイズスケジュールである．

\textbf{重要な性質}として，任意のステップ $t$ への分布が閉形式で得られる：
\begin{equation}
  q(\boldsymbol{x}_t | \boldsymbol{x}_0) = \mathcal{N}(\boldsymbol{x}_t; \sqrt{\bar{\alpha}_t}\boldsymbol{x}_0, (1-\bar{\alpha}_t)\mathbf{I}),
  \label{eq:forward-closed}
\end{equation}
ここで $\bar{\alpha}_t = \prod_{s=1}^t (1-\beta_s)$．

\subsubsection{逆向き過程}

生成は逆の過程で行う．$\boldsymbol{x}_T \sim \mathcal{N}(\mathbf{0}, \mathbf{I})$ から開始し，
段階的にノイズを除去する：
\begin{equation}
  p_\theta(\boldsymbol{x}_{t-1}|\boldsymbol{x}_t) = \mathcal{N}(\boldsymbol{x}_{t-1}; \boldsymbol{\mu}_\theta(\boldsymbol{x}_t, t), \Sigma_\theta(t)).
  \label{eq:reverse-process}
\end{equation}

\subsubsection{学習目標：ノイズ予測}

前向き過程で $\boldsymbol{x}_t = \sqrt{\bar{\alpha}_t}\boldsymbol{x}_0 + \sqrt{1-\bar{\alpha}_t}\boldsymbol{\epsilon}$
（$\boldsymbol{\epsilon} \sim \mathcal{N}(\mathbf{0}, \mathbf{I})$）とすると，
$\boldsymbol{\epsilon}$ を予測することで逆向き過程の平均を計算できる．

\textbf{損失関数}（Simplified ELBO）：
\begin{equation}
  \mathcal{L}_{\mathrm{simple}} = \mathbb{E}_{t, \boldsymbol{x}_0, \boldsymbol{\epsilon}}\left[\|\boldsymbol{\epsilon} - \boldsymbol{\epsilon}_\theta(\boldsymbol{x}_t, t)\|^2\right].
  \label{eq:ddpm-loss}
\end{equation}

このシンプルな目標関数で，非常に高品質な生成が実現される．

\subsection{拡散モデルと VAE の関係}

拡散モデルは\textbf{階層的VAE}として解釈できる．
前向き過程を変分分布，逆向き過程を生成分布と見なすと，
ELBOの形で統一的に扱える．

\subsection{条件付き生成}

\subsubsection{Classifier-Free Guidance}

条件 $\boldsymbol{c}$ が与えられたとき，条件付きと無条件のノイズ予測を混合する：
\begin{equation}
  \tilde{\boldsymbol{\epsilon}}_\theta = (1+w)\boldsymbol{\epsilon}_\theta(\boldsymbol{x}_t, t, \boldsymbol{c}) - w\boldsymbol{\epsilon}_\theta(\boldsymbol{x}_t, t, \emptyset),
  \label{eq:cfg}
\end{equation}
ここで $w$ はガイダンス強度である．

\subsubsection{クロスアテンションによる条件統合}

テキストなどの複雑な条件は，クロスアテンションで統合する：
\begin{equation}
  \mathrm{CrossAttn}(\mathbf{Q}, \mathbf{K}_c, \mathbf{V}_c) = \mathrm{softmax}\left(\frac{\mathbf{Q}\mathbf{K}_c^\top}{\sqrt{d}}\right)\mathbf{V}_c.
\end{equation}

\subsection{VAE と拡散モデルの比較}

表\ref{tab:vae-vs-diffusion}に両者の比較を示す．

\begin{table}[hbtp]
  \centering
  \caption{VAE と拡散モデルの比較．}
  \label{tab:vae-vs-diffusion}
  \begin{tabular}{lcc}
    \toprule
    特性 & VAE & 拡散モデル \\
    \midrule
    生成品質 & 中程度 & 優秀 \\
    推論速度 & 高速（1ステップ） & 遅い（多ステップ） \\
    訓練安定性 & 安定 & 比較的安定 \\
    潜在空間の獲得 & あり & なし \\
    \bottomrule
  \end{tabular}
\end{table}

% ------------------------------------------------------------------
\section{Physics-Informed Neural Networks（PINN）}
\label{sec:pinn}
% ------------------------------------------------------------------

\subsection{物理法則の組み込み}

\textbf{Physics-Informed Neural Networks（PINN）}\cite{Raissi2019}は，
物理法則（偏微分方程式など）をニューラルネットワークの学習に組み込む手法である．

交通流の連続方程式を例にとる：
\begin{equation}
  \frac{\partial \rho}{\partial t} + \frac{\partial (\rho v)}{\partial x} = 0,
  \label{eq:continuity}
\end{equation}
ここで $\rho$ は交通密度，$v$ は速度である．

\subsection{PINN の損失関数}

PINNの損失関数は，データフィッティング項と物理制約項から構成される：
\begin{equation}
  \mathcal{L} = \mathcal{L}_{\mathrm{data}} + \lambda \mathcal{L}_{\mathrm{physics}},
  \label{eq:pinn-loss}
\end{equation}
ここで：
\begin{align}
  \mathcal{L}_{\mathrm{data}} &= \frac{1}{N} \sum_{i=1}^{N} \|u_\theta(\boldsymbol{x}_i) - u_i\|^2, \\
  \mathcal{L}_{\mathrm{physics}} &= \frac{1}{M} \sum_{j=1}^{M} \|f_\theta(\boldsymbol{x}_j)\|^2.
\end{align}

$f_\theta$ は物理方程式の残差であり，自動微分により計算される．

\subsection{交通流への応用}

交通流モデルにPINNを適用することで，
物理的に整合性のある予測が可能になる．

\begin{example}
  LWR（Lighthill-Whitham-Richards）モデル：
  \begin{equation}
    \frac{\partial \rho}{\partial t} + \frac{\partial Q(\rho)}{\partial x} = 0,
  \end{equation}
  において，基本ダイアグラム $Q(\rho) = \rho v_f (1 - \rho/\rho_{\max})$ を
  PINNで学習する問題を考える．
  観測データと物理制約を組み合わせた損失関数を設計せよ．
\end{example}

\subsection{PINN の利点と課題}

\textbf{利点}：
\begin{itemize}
  \item 少量のデータでも物理的に整合性のある予測が可能
  \item 解釈可能性が高い
  \item 境界条件・初期条件の厳密な満足
\end{itemize}

\textbf{課題}：
\begin{itemize}
  \item 複雑な物理系での収束の困難さ
  \item 物理パラメータの推定における不確実性
  \item 計算コスト（自動微分のオーバーヘッド）
\end{itemize}

% ------------------------------------------------------------------
\section{交通・都市工学への応用}
\label{sec:generative-traffic}
% ------------------------------------------------------------------

\subsection{軌跡生成}

都市計画やシミュレーションでは，現実的な交通軌跡を合成する必要がある．

\textbf{VAEベースの軌跡生成}：
\begin{enumerate}
  \item エンコーダ：軌跡を入力し，意図を圧縮
  \item デコーダ：潜在意図から軌跡を再構成
  \item 生成：新しい $\boldsymbol{z}$ をサンプリング → 新規軌跡を生成
\end{enumerate}

\textbf{拡散モデルベースの軌跡生成}：
\begin{enumerate}
  \item 実軌跡にノイズを段階的に追加
  \item ノイズから軌跡を段階的に除去
  \item より自然で多様な軌跡が生成可能
\end{enumerate}

\subsection{条件付き都市計画}

「駅周辺の歩行者流量が高い地区」という条件で，
建築配置パターンを生成する応用が考えられる．

\begin{enumerate}
  \item 条件付き拡散モデルを訓練
  \item Classifier-Free Guidanceで生成品質を調整
  \item 設計者が多数の候補から最適案を選択
\end{enumerate}

\subsection{交通シミュレーションの高速化}

PINNを用いることで，交通流シミュレーションを
数値解法より高速に実行できる可能性がある．

\begin{column}{サロゲートモデルとしてのPINN}
  従来の数値シミュレーション（有限差分法など）は計算コストが高い．
  PINNを「サロゲートモデル」として訓練すれば，
  一度学習した後は高速な予測が可能になる．
  これは交通シナリオの多数評価が必要な最適化問題で特に有用である．
\end{column}

% ------------------------------------------------------------------
\section{まとめと今後の展望}
\label{sec:generative-summary}
% ------------------------------------------------------------------

本章では，生成モデルの統一的な理解を目指した．

\textbf{学んだこと}：
\begin{enumerate}
  \item \textbf{潜在変数モデル}：データの低次元構造を捉える
  \item \textbf{変分推論}：計算困難な積分をELBOで下限から制御
  \item \textbf{VAE}：安定で解釈可能な実装
  \item \textbf{拡散モデル}：現在の生成品質の最前線
  \item \textbf{PINN}：物理法則を組み込んだ学習
\end{enumerate}

\textbf{今後の展望}：
\begin{itemize}
  \item \textbf{スケーラビリティ}：大規模・高次元データへの対応
  \item \textbf{条件付き生成の洗練}：テキスト，マルチモーダル入力への対応
  \item \textbf{計算効率}：拡散モデルの多ステップ生成を削減する研究
  \item \textbf{応用拡大}：科学計算，分子設計，構造設計など新領域への展開
\end{itemize}

生成モデルは，単なる画像生成ツールではなく，
\textbf{データの本質的な構造を学習する深い数学的フレームワーク}である．
